\documentclass[ekobook.tex]{subfiles}

\begin{document}\setWPP
\setlist[itemize]{
  itemsep=0pt,
  parsep=0pt,      % space between paragraph lines inside one item
  topsep=0.5ex,    % space before/after the whole list (usually keep a little)
  partopsep=0pt
}
\lsection{Ekonomické pojmy v praxi}
\qaw{Vysvětli pojmy makroekonomie a mikroekonomie}{%
\textbl{Makroekonomie} -- zkoumá ekonomiku jako \texthl{celek} (inflace, nezaměstnanost, HDP, hospodářský růst, státní rozpočet).\\
\textbl{Mikroekonomie} -- zabývá se chováním \texthl{jednotlivců} a malých ekonomických subjektů (domácnosti, firmy, spotřebitelé, ceny jednotlivých statků).
}
\begin{multicols}{2}
\qaw{Vysvětli:}{%
\textbl{zvykový systém}
-- nejstarší systém založený na \texthl{tradici} a zvyklostech dané země/společnosti\\
\textbl{příkazový systém}
-- vláda nebo centrální orgán \texthl{závazně určuje}, co, kolik a pro koho se má vyrábět\\
\textbl{tržní systém}
-- ekonomika řízená \texthl{nabídkou a poptávkou} s minimálním zásahem státu\\
\textbl{smíšený systém}
-- kombinuje prvky tržní a příkazové ekonomiky; stát stanovuje mantinely, ve kterých funguje volný trh\\ (většina současných ekonomik)
}
\columnbreak
\qaw{Co je to potřeba}{%
\begin{itemize}[label={\color{WPP}\faSlackHash}]
\item \textbl{hmotná} -- jídlo, oblečení, byt
\item \textbl{nehmotná} -- vzdělání, znalosti, dovednosti
\item \textbl{nezbytná} -- jídlo, bydlení, základní ošacení
\item \textbl{zbytná} -- dovolená v zahraničí, luxusní auto, paměťová karta
\item \textbl{fyzická} -- jídlo, pití, teplo
\item \textbl{duševní} -- četba knihy, poslech hudby, přátelství
\item \textbl{současná} -- okamžitá potřeba (hlad, žízeň)
\item \textbl{budoucí} -- spoření na důchod, vzdělání dětí
\end{itemize}
}
\end{multicols}
\wbreak
\qaw{Složky životní úrovně}{%
\begin{itemize}[label=\textcolor{WPP}{$\hookrightarrow$}]
\item osobní spotřeba obyvatelstva
\item společenská spotřeba (veřejné služby)
\item {úroveň bydlení}
\item volný čas a rekreace
\item sociální jistoty a zdravotní péče
\item pracovní podmínky
\item kvalita {životního prostředí}
\end{itemize}
}
\qaw{Uveď příklad státu s vysokou životní úrovní a s nízkou životní úrovní}{%
\textbl{Vysoká životní úroveň:} Švýcarsko, Norsko, Austrálie, Dánsko, Island\\
\textbl{Nízká životní úroveň:} některé země subsaharské Afriky, Afghánistán, Jemen (v Evropě např. Moldavsko, Ukrajina)
}
\qaw{Vyjmenuj výrobní faktory}{%
\textbl{práce} -- lidský faktor (fyzická i duševní práce)\\
\textbl{přírodní zdroje} -- půda, nerosty, voda, lesy, klima\\
\textbl{kapitál} -- stroje, budovy, zařízení, peněžní prostředky (vytvořené výrobní prostředky)
}
\qaw{Co je nabídka, poptávka}{%
\textbl{Nabídka} -- množství zboží a služeb, které jsou výrobci ochotni dodat na trh při dané ceně (\texthl{čím vyšší cena, tím vyšší nabídka}).\\
\textbl{Poptávka} -- množství zboží a služeb, které jsou spotřebitelé ochotni a schopni koupit při dané ceně (\texthl{čím nižší cena, tím vyšší poptávka}).
}
\qaw{Substitut, komplement}{%
\textbl{Substitut} -- zboží, které může nahradit jiné zboží (kuřecí maso × vepřové maso, jablka × hrušky).\\
\textbl{Komplement} -- zboží, které se spotřebovává společně (auto a benzín, pračka a prací prášek).
}
\vspace*{1em}
\wbreak
\qaw{RPSN}{%
\texthl{Roční procentní sazba nákladů} -- ukazuje \texthl{skutečné roční náklady úvěru} v procentech. Zahrnuje úrok i všechny poplatky (za sjednání, vedení účtu, pojištění apod.).\\
\wemph{Čím nižší RPSN, tím výhodnější úvěr.}
}
\qaw{Plátce daně X poplatník}{%
\textbl{Plátce daně} -- osoba, která daň \texthl{vypočítává, vybírá a odvádí} státu (např. zaměstnavatel u daně z příjmů, prodejce u DPH).\\
\textbl{Poplatník} -- osoba, \texthl{ze které je daň reálně hrazena} (zaměstnanec, konečný spotřebitel).
}
\qaw{Zdaňovací období, sazba daně, sleva na dani}{%
\textbl{Zdaňovací období} -- období, za které se daň počítá a platí (obvykle \texthl{rok} nebo \texthl{měsíc}).\\
\textbl{Sazba daně} -- procentuální vyjádření daně (např. daň z příjmů 15 \%, DPH 21 \%).\\
\textbl{Sleva na dani} -- snižuje vypočtenou daň (sleva na poplatníka, na dítě, na manžela/manželku, ZTP).
}
\qaw{Daňový bonus, přeplatek a nedoplatek na dani}{%
\textbl{Daňový bonus} -- vzniká, když součet slev (zejména na děti) převýší vypočtenou daň; stát rozdíl \texthl{vyplatí}.\\
\textbl{Přeplatek na dani} -- zaplatili jsme více, než byla daňová povinnost; stát přeplatek \texthl{vrátí}.\\
\textbl{Nedoplatek na dani} -- zaplatili jsme méně, než byla daňová povinnost; musíme \texthl{doplatit}.
}
\qaw{Opotřebení majetku}{%
Ztráta hodnoty majetku v průběhu času. Dělí se na:\\
\textbl{fyzické opotřebení} -- opotřebení používáním (auto, stroj)\\
\textbl{morální opotřebení} -- technické nebo morální zastarávání (starý počítač, počítačový program)
}
\qaw{Nominální a reálná mzda}{%
\textbl{Nominální mzda} -- peněžní částka uvedená na výplatní pásce.\\
\textbl{Reálná mzda} -- skutečná kupní síla mzdy; co si za ni zaměstnanec může koupit.
}
\qaw{Inflace, deflace a spotřební koš}{%
\textbl{Inflace} -- růst \texthl{cenové hladiny} v ekonomice.\\
\textbl{Deflace} -- pokles \texthl{cenové hladiny} (obvykle nevhodný jev).\\
\textbl{Spotřební koš} -- soubor vybraných výrobků a služeb (přibližně 700 položek), podle kterého se sleduje vývoj cen a inflace.}
\qaw{Kontokorentní úvěr}{%
Úvěr, který umožňuje \texthl{čerpat peníze z běžného účtu do minusu} (povolené přečerpání).
Má obvykle \texthl{velmi vysoký úrok}.
}
\qaw{Hypoteční úvěr}{%
Dlouhodobý úvěr určený především na \texthl{nákup, výstavbu nebo rekonstrukci nemovitosti}.
Zajištěn zástavním právem k nemovitosti.
}
\qaw{Lombardní úvěr}{%
Úvěr zajištěný \texthl{movitou věcí} (cenné papíry, zboží, šperky).
Rizikovější než hypotéka, proto vyšší úrok.
}
\qaw{Debetní X kreditní karta}{%
\textbl{Debetní karta} -- platíme z vlastních peněz na bankovním účtu.\\
\textbl{Kreditní karta} -- umožňuje \texthl{čerpat úvěr} banky (platby do minusu), obvykle s bezúročným obdobím.
}
\qaw{Spotřební úvěr}{%
Úvěr poskytovaný spotřebiteli. Může být \texthl{účelový} (na konkrétní věc, např. auto) nebo \texthl{neúčelový} (na cokoli).
}
\wpage
\end{document}